\documentclass[UKenglish,10pt,a4paper]{report}
\setlength\textwidth{145mm}
\setlength\textheight{247mm}
\setlength\oddsidemargin{5mm}
\setlength\evensidemargin{15mm}
\setlength\topmargin{0mm}
\setlength\headsep{0mm}
\setlength\headheight{0mm}
% \openright zařídí, aby následující text začínal na pravé straně knihy
\let\openright=\clearpage
\renewcommand{\baselinestretch}{1.0}%řádkování
\usepackage{booktabs}
\usepackage{multirow}
\usepackage{tabu}
\usepackage[hyphens]{url}
\usepackage{color}
\usepackage{listings}
\usepackage[utf8]{inputenc}
\usepackage{csquotes}
\usepackage[backend=biber]{biblatex}
\usepackage[pdftex]{graphicx}
\usepackage{epstopdf}
\usepackage{wrapfig}
\usepackage[section]{placeins}
\usepackage{amsfonts}
\usepackage{subcaption}
\usepackage{chronosys}
\usepackage{textcomp}
\usepackage{framed}
\usepackage{float}
\usepackage[czech, UKenglish]{babel}
\usepackage{multicol}
\usepackage{breakurl}
\usepackage[left]{eurosym}
\usepackage[binary-units]{siunitx}
\usepackage{makecell}
\usepackage{isodate}





\newcommand\ytl[2]{
\parbox[b]{8em}{\hfill{\color{cyan}\bfseries\sffamily #1}~$\cdots\cdots$~}\makebox[0pt][c]{$\bullet$}\vrule\quad \parbox[c]{7.5cm}{\vspace{7pt}\color{red!40!black!80}\raggedright\sffamily #2\\[7pt]}\\[-3pt]}

\usepackage[T1]{fontenc}
\usepackage{amsmath}
%\usepackage[toc,page]{appendix}
\usepackage[acronym,toc]{glossaries} 

\usepackage[unicode,breaklinks]{hyperref}  % Musí být za všemi ostatními balíčky
\usepackage{bookmark}




%\usepackage{pdfpages}

\addbibresource{ref.bib}


%% Pokud tiskneme oboustranně:
% \documentclass[12pt,a4paper,twoside,openright]{report}
% \setlength\textwidth{145mm}
% \setlength\textheight{247mm}
% \setlength\oddsidemargin{15mm}
% \setlength\evensidemargin{0mm}
% \setlength\topmargin{0mm}
% \setlength\headsep{0mm}
% \setlength\headheight{0mm}
% \let\openright=\cleardoublepage

%% Pokud používáte csLaTeX (doporučeno):
%\usepackage[czech]{babel}
%% Pokud nikoliv:

\begingroup
    \makeatletter
    \catcode`\-=\active
    \AtBeginDocument{
    \def\@@@cmidrule[#1-#2]#3#4{\global\@cmidla#1\relax
        \global\advance\@cmidla\m@ne
        \ifnum\@cmidla>0\global\let\@gtempa\@cmidrulea\else
        \global\let\@gtempa\@cmidruleb\fi
        \global\@cmidlb#2\relax
        \global\advance\@cmidlb-\@cmidla
        \global\@thisrulewidth=#3
        \@setrulekerning{#4}
        \ifnum\@lastruleclass=\z@\vskip \aboverulesep\fi
        \ifnum0=`{\fi}\@gtempa
        \noalign{\ifnum0=`}\fi\futurenonspacelet\@tempa\@xcmidrule}
    }
\endgroup


%% Použité kódování znaků: obvykle latin2, cp1250 nebo utf8:
%\usepackage[cp1250]{inputenc}

%\usepackage[T1]{fontenc} %tohle by mělo udělat Times New Roman 
%% Ostatní balíčky

%\usepackage{mathtools}
%\usepackage{amsthm}

%% Balíček hyperref, kterým jdou vyrábět klikací odkazy v PDF,
%% ale hlavně ho používáme k uložení metadat do PDF (včetně obsahu).
%% POZOR, nezapomeňte vyplnit jméno práce a autora.
\hypersetup{pdftitle=Optical communications}
\hypersetup{pdfauthor=Bc. et Bc. Jaroslav Svoboda, pdfkeywords={optical,network}}
%\hypersetup{hyperfootnotes=true}
\makeglossaries
%%% Drobné úpravy stylu

% Tato makra přesvědčují mírně ošklivým trikem LaTeX, aby hlavičky kapitol
% sázel příčetněji a nevynechával nad nimi spoustu místa. Směle ignorujte.
\makeatletter
\def\@makechapterhead#1{
  {\parindent \z@ \raggedright \normalfont
   \Huge\bfseries \thechapter. #1
   \par\nobreak
   \vskip 20\p@
}}
\def\@makeschapterhead#1{
  {\parindent \z@ \raggedright \normalfont
   \Huge\bfseries #1
   \par\nobreak
   \vskip 20\p@
}}
\makeatother

% Toto makro definuje kapitolu, která není očíslovaná, ale je uvedena v obsahu.
\def\chapwithtoc#1{
\chapter*{#1}}
\renewcommand{\familydefault}{\rmdefault}
\begin{document}
\selectlanguage{UKenglish}
% Trochu volnější nastavení dělení slov, než je default.
\lefthyphenmin=2
\righthyphenmin=2
\sisetup{range-phrase=--}
%%% Titulní strana práce
\begin{titlepage}


\pagestyle{empty}
\begin{center}

\large

TECHNICAL UNIVERSITY OF MADRID

\medskip

SCHOOL OF TELECOMMUNICATIONS SYSTEMS AND ENGINEERING

\vfill

{\bf\Large Semester Project}

\vfill
\centerline{\mbox{\includegraphics[width=60mm]{img/logo2.eps}\qquad{\includegraphics[width=60mm]{img/logo.eps}}}}
			
\vfill
\vspace{5mm}

{\LARGE Bc. et Bc. Jaroslav Svoboda}

{\LARGE Jaime Sancho}

{\LARGE Diallo Elhadj Sadou}

\vspace{10mm}

% Název práce přesně podle zadání
{\LARGE\bfseries } Smart Lighting

\vfill

% Název katedry nebo ústavu, kde byla práce oficiálně zadána
% (dle Organizační struktury MFF UK)
UBIQUITOUS AND SECURE NETWORKS AND SERVICES

\vfill

\begin{tabular}{rl}

Lecturer: & Jesus Rodriguez Molina  \\
\noalign{\vspace{2mm}}
Subject coordinator: & Jose Fernan Martinez Ortega \\
\noalign{\vspace{2mm}}
Degree programme: & Master in Systems and Services Engineering for the Information society \\
\end{tabular}

\vfill

% Zde doplňte rok
Madrid 2018

\end{center}
\end{titlepage}
\newpage

%%% Následuje vevázaný list -- kopie podepsaného "Zadání bakalářské práce".
%%% Toto zadání NENÍ součástí elektronické verze práce, nescanovat.


\openright

\noindent



%%% Strana s automaticky generovaným obsahem bakalářské práce. U matematických
%%% prací je přípustné, aby seznam tabulek a zkratek, existují-li, byl umístěn
%%% na začátku práce, místo na jejím konci.



\openright
\pagestyle{plain}
\setcounter{page}{1}
\tableofcontents



\chapter{Introduction and Executive Summary}
\chapter{Scheduling of workload}
\chapter{Requirement Analysis}
\section{General use}
\begin{figure}[ht]
	\caption{General use}
	\label{fig:general_use}
	\vspace{-1em}
	\centering
	\includegraphics[width=0.48\textwidth]{{"./img/Drawio/Use Cases/General use case diagram"}.png}
\end{figure}

\begin{table}[!ht]
	\centering
	\caption{General use}
	\label{tab:general_use}
	\begin{tabu}{|l|X[l]|}
		\hline
		Scope and Objectives & \noindent\parbox[c]{\hsize}{The lights on the lampposts switch on/off depending on the information of the movement and light sensor of the motes (if the light is low and the movement sensor detects something, the lights turn on, otherwise they turn off). The data collection centre receives periodically this data and it is capable to change some parameters or request some information of the motes.} \\ \hline
		Actors               & \noindent\parbox[c]{\hsize}{\begin{itemize}
								\item Pedestrians/vehicles
								\item Data control department of the city
								\end{itemize} } \\ \hline
		Preconditions        & \noindent\parbox[c]{\hsize}{\begin{itemize}
								\item The sensor nodes are on
								\item A pedestrian/vehicle comes/goes away from the node
								\end{itemize} }  \\ \hline
		Post-conditions      & \noindent\parbox[c]{\hsize}{ \begin{itemize}
								\item The lights will turn on/off
								\item All the measured information will be sent to the city data manager
								\end{itemize}} \\ \hline
		Sequence Description & \noindent\parbox[c]{\hsize}{\begin{enumerate}
								\item Obtain measurements of movement, temperature, humidity and light.
								\item If there is a change in the movement sensor and the light measured is low, toggle the lights on the node.
								\item Send the status of light to the data control department of the city.
								\end{enumerate}} \\ \hline
		Exceptions           & \noindent\parbox[c]{\hsize}{The node is not able to communicate with the data control department of the city.} \\ \hline
	\end{tabu}

\end{table}

\newpage
\section{Receive Periodic Data}
\begin{figure}[!ht]
	\caption{Receive Periodic Data}
	\label{fig:peridic_data}
	\vspace{-1em}
	\centering
	\includegraphics[width=0.78\textwidth]{{"./img/Drawio/Use Cases/Periodically receive data case diagram"}.png}
\end{figure}

\begin{table}[ht]
	\centering
	\caption{Receive Periodic Data}
	\label{tab:periodic_data}
	\begin{tabu}{|l|X[l]|}
		\hline
		Scope and Objectives & \parbox[c]{\hsize}{The motes are placed in the lampposts of the street getting information of the movement and the light. This information added to the LED status of the lamppost is sent to the data collection centre periodically.} \\ \hline
		Actors               & \noindent\parbox[c]{\hsize}{Data control department of the city}      \\ \hline
		Preconditions        & \noindent\parbox[c]{\hsize}{	\begin{itemize}
								\item The sensor nodes are on
								\item The sensor nodes have sufficient memory to store measured values when connection is off
								\end{itemize}}  \\ \hline
		Post-conditions      & \noindent\parbox[c]{\hsize}{\begin{itemize}
								\item The measured value is stored in memory
								\item The measured value is sent to the data control department of the city 
								\end{itemize} } \\ \hline
		Sequence Description & \noindent\parbox[c]{\hsize}{ \begin{enumerate}
								\item Obtain measurement
								\item Check the network
								\item If network not available store data
								\item Send all data to the data control department of the city 
								\end{enumerate} } \\ \hline
		Exceptions           & \noindent\parbox[c]{\hsize}{In case of error, report to the data control department of the city} \\ \hline
	\end{tabu}
\end{table}

\newpage

\section{Request data}
\begin{figure}[!ht]
	\caption{Request Data}
	\label{fig:request_data}
	\vspace{-1em}
	\centering
	\includegraphics[width=0.68\textwidth]{{"./img/Drawio/Use Cases/Request information case diagram"}.png}
\end{figure}

\begin{table}[!ht]
	\centering
	\caption{Request Data}
	\label{tab:request_data}
	\begin{tabu}{|l|X[l]|}
		\hline
		Scope and Objectives & \noindent\parbox[c]{\hsize}{The data collection centre have the possibility to request the information of the light and LED status of the motes and also the parameters that are using the motes. Furthermore, a ping functionality exists to check if the mote is alive.} \\ \hline
		Actors               & \noindent\parbox[c]{\hsize}{Data control department of the city}  \\ \hline
		Preconditions        & \noindent\parbox[c]{\hsize}{\begin{itemize}
				 \item The sensor nodes are on
				 \item The sensor node is online to get the request
			 \end{itemize} }   \\ \hline
		Post-conditions      & \noindent\parbox[c]{\hsize}{The measured value and statuses are sent to the data control department of the city} \\ \hline
		Sequence Description & \noindent\parbox[c]{\hsize}{\begin{enumerate}
				\item Check if node is online
				\item If not then repeat x times
				\item Request current data from node
			\end{enumerate}} \\ \hline
		Exceptions           & \noindent\parbox[c]{\hsize}{The sensor node does not respond after x requests}  \\ \hline
	\end{tabu}
\end{table}

\newpage

\section{Change parameters}
\begin{figure}[!ht]
	\caption{Change parameters}
	\label{fig:change_parameters}
	\vspace{-1em}
	\centering
	\includegraphics[width=0.78\textwidth]{{"./img/Drawio/Use Cases/Change parameters case diagram"}.png}
\end{figure}

\begin{table}[!ht]
	\centering
	\caption{Change parameters}
	\label{tab:change_parameters}
	\begin{tabu}{|l|X[l]|}
		\hline
		Scope and Objectives & \noindent\parbox[c]{\hsize}{The data collection center is able to change the threshold of the light sensor and also has the possibility of switch on/off the lights in case of emergency.}     \\ \hline
		Actors               & \noindent\parbox[c]{\hsize}{Data control department of the city}                                                                                                                               \\ \hline
		Preconditions        & \noindent\parbox[c]{\hsize}{\begin{itemize}\item The sensor nodes are on \item The sensor node is online to get the request \item Led light is working\end{itemize} }  \\ \hline
		Post-conditions      & \noindent\parbox[c]{\hsize}{\begin{itemize}\item The light is turned on/off \item Parameters are modified \item The status is sent to the data control department of the city\end{itemize}}  \\ \hline
		Sequence Description & \noindent\parbox[c]{\hsize}{\begin{enumerate}\item Check if node is online \item If not then repeat x times \item Request switching on/off the light or change of parameters \item Send ack when light turns on or parameters are changed \end{enumerate}} \\ \hline
		Exceptions           & \noindent\parbox[c]{\hsize}{\begin{itemize}\item The sensor node does not respond after x requests \item The light does not turn on/off\end{itemize}}         \\ \hline
	\end{tabu}
\end{table}

\chapter{Project Design}
\section{UML Component Diagrams}

\begin{figure}[!ht]
	\caption{Movement sensor}
	\label{fig:movement_sensor_component}
	\vspace{-1em}
	\centering
	\includegraphics[width=0.78\textwidth]{{"./img/Drawio/Components/Normal sensor"}.png}
\end{figure}
\begin{figure}[!ht]
	\caption{Light sensor}
	\label{fig:light_sensor_component}
	\vspace{-1em}
	\centering
	\includegraphics[width=0.78\textwidth]{{"./img/Drawio/Components/Light Sensor"}.png}
\end{figure}
\begin{figure}[!ht]
	\caption{Data collection centre}
	\label{fig:data_collection_centre_component}
	\vspace{-1em}
	\centering
	\includegraphics[width=0.78\textwidth]{{"./img/Drawio/Components/datacollection"}.png}
\end{figure}
\begin{figure}[!ht]
	\caption{Aggregator}
	\label{fig:aggregator_component}
	\vspace{-1em}
	\centering
	\includegraphics[width=0.78\textwidth]{{"./img/Drawio/Components/Aggregator"}.png}
\end{figure}

\section{UML Class Diagrams}
\begin{figure}[!ht]
	\caption{Movement sensor}
	\label{fig:movement_sensor_class}
	\vspace{-1em}
	\centering
	\includegraphics[width=0.78\textwidth]{{"./img/Drawio/Class Diagrams/NormalSensor"}.png}
\end{figure}
\begin{figure}[!ht]
	\caption{Light sensor}
	\label{fig:light_sensor_class}
	\vspace{-1em}
	\centering
	\includegraphics[width=0.78\textwidth]{{"./img/Drawio/Class Diagrams/LightSensor"}.png}
\end{figure}
\begin{figure}[!ht]
	\caption{Sink}
	\label{fig:sink_class}
	\vspace{-1em}
	\centering
	\includegraphics[width=0.78\textwidth]{{"./img/Drawio/Class Diagrams/sink"}.png}
\end{figure}
\begin{figure}[!ht]
	\caption{Aggregator}
	\label{fig:aggregator_class}
	\vspace{-1em}
	\centering
	\includegraphics[width=0.78\textwidth]{{"./img/Drawio/Class Diagrams/Aggregator"}.png}
\end{figure}

\section{Sequence Diagrams}
\begin{figure}[!ht]
	\caption{Light Broadcast}
	\label{fig:light_broadcast_sequence}
	\vspace{-1em}
	\centering
	\includegraphics[width=0.98\textwidth]{{"./img/Drawio/Sequence Diagrams/LightBroadcast"}.png}
\end{figure}
\renewcommand{\acronymname}{List of acronyms}
\chapter{Testing}
\chapter{Conclusion}


\printglossaries



%\bibliographystyle{plain} 
\printbibliography
%%% Přílohy k bakalářské práci, existují-li (různé dodatky jako výpisy programů,
%%% diagramy apod.). Každá příloha musí být alespoň jednou odkazována z vlastního
%%% textu práce. Přílohy se číslují.
\openright
\end{document}